%% sections/06-procedure.tex
%% Step-by-Step Procedure
%% IEC/IEEE 82079-1:2019 — Instructional Information (§5)
%% ─────────────────────────────────────────────────────────

\section{Procedure}
\label{sec:procedure}

%%
%% WRITING GUIDE — PROCEDURE SECTION:
%%
%%   Structure rules:
%%   - Use \begin{steplist} ... \end{steplist} for numbered steps.
%%   - End every major step group with \outcome{...}.
%%   - Put WARNING / CAUTION / NOTE boxes immediately BEFORE the step
%%     that needs them, not after.
%%   - Embed code with \begin{pythoncode}[title] or \begin{acscode}[title].
%%   - Use \code{text} for inline parameter/command names.
%%   - Link to figures and tables with \cref{label}.
%%

\subsection{Phase 1 — Enable Servo and Zero Sensor}
\label{sec:proc:phase1}

\begin{warningbox}
  Ensure no workpiece or tooling is below the Z-axis head before enabling
  the servo.  The head will hold its current position when the servo
  enables.
\end{warningbox}

\begin{steplist}

  \step Enable the Z-axis servo from the ACS MMI:
        navigate to \menu{Axes > Control} and click \textbf{Enable} for
        Axis~0 (Z-axis), or run the following ACSPL+ snippet in the
        SPiiPlus MMI Terminal:

        \begin{acscode}[Enable Z-axis servo (Axis 0)]
\begin{lstlisting}
; Enable servo on Axis 0
ENABLE 0
WAIT 100        ; Wait 100 ms for servo to stabilise
PRINT "Servo state: ", MFLAGS(0)
\end{lstlisting}
        \end{acscode}

  \outcome{The axis MFLAGS bit~0 (MSF\_ENABLED) is set.
           The ACS MMI axis indicator turns \textcolor{TipGreen}{\textbf{green}}.}

  \step Zero the force sensor reading.
        Apply \textbf{no load} to the load cell (free-hanging Z-axis),
        then run:

        \begin{acscode}[Zero the force sensor (remove offset)]
\begin{lstlisting}
; Read current AIN voltage and store as offset
GLOBAL REAL force_offset
force_offset = AIN(0, 1)   ; AIN channel 1
PRINT "Force offset set to: ", force_offset, " V"
\end{lstlisting}
        \end{acscode}

  \outcome{\code{force\_offset} is stored.  The value should be within
           ±\SI{0.05}{\volt} of zero; a larger value suggests sensor
           misalignment or electrical noise — see \cref{sec:troubleshooting}.}

\end{steplist}


\subsection{Phase 2 — Configure the Force Control Loop}
\label{sec:proc:phase2}

\begin{steplist}

  \step Set the force sensor scale factor and PID gains in the ACS
        parameter table.  Use the baseline values from \cref{tab:pid-defaults}
        and update \acsparam{FORCE\_SCALE} to match your sensor sensitivity:

        \begin{acscode}[Set force loop parameters]
\begin{lstlisting}
; Force sensor scale: Newton per Volt
; Adjust SENSOR_SCALE to your calibrated value
GLOBAL REAL SENSOR_SCALE
SENSOR_SCALE = 100.0        ; [N/V] — Kistler 9317C default

; PID gains (see Section 3.3 for background)
XVEL(0) = 0.0               ; Velocity limit while tuning
KP(0) = 2.5
KI(0) = 0.08
KD(0) = 0.001
KP_FFW(0) = 0.0
PRINT "Gains set. KP=", KP(0), " KI=", KI(0)
\end{lstlisting}
        \end{acscode}

  \step Set the force setpoint and safety current limit:

        \begin{acscode}[Set target force and current limit]
\begin{lstlisting}
; Start with a LOW force for first test
GLOBAL REAL force_setpoint
force_setpoint = 5.0        ; [N] — increase after verification

; Current limit (safety ceiling — NEVER exceed motor spec)
DCURR(0) = 1.5              ; [A] — max motor continuous current
ECURR(0) = 2.0              ; [A] — emergency current limit
\end{lstlisting}
        \end{acscode}

  \begin{cautionbox}
    Set \acsparam{force\_setpoint} to no more than \SI{5}{\newton} for
    first-time commissioning on a new machine or a new workpiece type.
    Increase in \SI{5}{\newton} steps only after confirming stable response.
  \end{cautionbox}

  \step Start the force control loop using the PBA force control ACSPL+
        program:

        \begin{acscode}[Start force control loop — ACSPL+ program]
\begin{lstlisting}
; Load the force control program from buffer 5
START 5
WAIT 200                    ; Wait for program to initialise

; Monitor force error for 1 second
GLOBAL INT i
FOR i = 1 TO 100
    PRINT "F_meas=", (AIN(0,1) - force_offset) * SENSOR_SCALE, " N"
    SLEEP 10
NEXT i
\end{lstlisting}
        \end{acscode}

\end{steplist}


\subsection{Phase 3 — Python Force Logging}
\label{sec:proc:phase3}

Run the Python logging script from the host PC to record force data during
a test cycle.  The script is provided in \cref{app:code-python}.

\begin{steplist}

  \step Open a terminal on the host PC and navigate to the project folder:

        \begin{pythoncode}[Launch force logger from terminal]
\begin{lstlisting}
# Navigate to the project folder
cd C:\PBA\ForceControl\scripts

# Run the logger (records for 30 seconds by default)
python force_logger.py --ip 10.0.0.100 --duration 30 --output force_log.csv
\end{lstlisting}
        \end{pythoncode}

  \step While the logger runs, gently lower the Z-axis until contact is
        made.  Observe the force value in the terminal output.

  \step After the logging duration completes, check the output file:

        \begin{pythoncode}[Verify log file]
\begin{lstlisting}
import pandas as pd

df = pd.read_csv("force_log.csv")
print(df.describe())
print(f"Max force recorded: {df['force_N'].max():.2f} N")
print(f"Std deviation:      {df['force_N'].std():.3f} N")
\end{lstlisting}
        \end{pythoncode}

\end{steplist}

\outcome{The CSV file contains time-stamped force readings.  The standard
         deviation of the force signal at steady state should be less than
         \SI{0.1}{\newton} RMS.  If not, see \cref{sec:troubleshooting}.}

\videolink{https://your-video-url-here.com}{Watch: Force ramp test — sensor feedback and ACS MMI display (45\,s)}
